\subsubsection{Each composition notices at most two folds}\label{ov:twocolors}

We now show that any composition $X;Y$ of color-determined relations
notices at most two folds (Lemma~\ref{lem:local}). The two-witness
property first rules out noticing at pairs away from the hub.

\paragraph{Away from the hub (Figure~\ref{fig:notice}).}
At pairs $(u,v)$ with $u,v\ne\hub$, composition agrees with
prediction.
Figure~\ref{fig:notice} uses the same coloring, operands $X,Y$, and fold
as Figure~\ref{fig:prediction}, but observes $e'=(u,v_\mu)$.

\begin{figure}[H]
  \centering
  \begin{tikzpicture}[
    >=Latex,
    every node/.style={font=\scriptsize},
    vtx/.style={circle,draw,fill=white,inner sep=1.5pt},
    edge/.style={line width=.55pt},
    recolored/.style={red!75!black,densely dashed,line width=.85pt},
    surviving/.style={line width=1.15pt},
    read/.style={densely dotted,line width=.55pt},
  ]
    % Same four vertices and positions before and after folding.
    % u is the vertex z of Figure 2, and w is its alternative witness.
    \begin{scope}
      \coordinate (u0) at (0,0);
      \coordinate (v0) at (2.65,0);
      \coordinate (h0) at (1.325,.98);
      \coordinate (w0) at (1.325,-.98);
      \draw[read] (u0)--node[below,pos=.57] {$\mu$} (v0);
      \draw[edge] (u0)--node[above left] {$\lambda$} (h0);
      \draw[edge] (h0)--node[above right] {$\mu$} (v0);
      \draw[surviving] (u0)--node[below left] {$\lambda$} (w0);
      \draw[surviving] (w0)--node[below right] {$\mu$} (v0);
      \node[vtx,label=left:{$u$}] at (u0) {};
      \node[vtx,label=right:{$v_\mu$}] at (v0) {};
      \node[vtx,label=above:{$\hub$}] at (h0) {};
      \node[vtx,label=below:{$w$}] at (w0) {};
      \node[font=\small] at (.1, 1.2) {$\mathcal M_0$};
    \end{scope}
    \draw[-{Latex[length=2mm]},line width=.7pt]
      (3.30,0)--node[above] {fold $\lambda$} (4.65,0);
    \begin{scope}[xshift=5.25cm]
      \coordinate (u1) at (0,0);
      \coordinate (v1) at (2.65,0);
      \coordinate (h1) at (1.325,.98);
      \coordinate (w1) at (1.325,-.98);
      \draw[read] (u1)--node[below,pos=.57] {$\mu$} (v1);
      \draw[recolored] (u1)--node[above left,black] {$\lambda^c$} (h1);
      \draw[edge] (h1)--node[above right] {$\mu$} (v1);
      \draw[surviving] (u1)--node[below left] {$\lambda$} (w1);
      \draw[surviving] (w1)--node[below right] {$\mu$} (v1);
      \node[vtx,label=left:{$u$}] at (u1) {};
      \node[vtx,label=right:{$v_\mu$}] at (v1) {};
      \node[vtx,label=above:{$\hub$}] at (h1) {};
      \node[vtx,label=below:{$w$}] at (w1) {};
      \node[font=\small] at (.1, 1.2) {$\mathcal M_\lambda$};
    \end{scope}
  \end{tikzpicture}
  \caption{Away from the hub, the fold invalidates the path through $\hub$
    but leaves the second witness path through $w$ unchanged.
    Only these two witness paths are shown.}
  \label{fig:notice}
\end{figure}

In $\mathcal M_0$, both paths have colors $(\lambda,\mu)$.
After the fold, the path through $\hub$ is invalidated but the path through $w$ survives.
Hence:
$(X^{(\lambda)};Y^{(\lambda)})(e')
=1
=(X;Y)^{(\lambda)}(e')$.

More generally, between any distinct $u,v\ne\hub$, the two-witness
property preserves every two-color path pattern via a witness other
than $\hub$. Loops away from $\hub$ retain edges of every color, and
endpoint witnesses are unchanged. Thus composition cannot violate
prediction away from the hub.

\paragraph{At the hub (Figure~\ref{fig:prediction} revisited).}

Only pairs involving $\hub$ remain.
In Figure~\ref{fig:prediction}, $X$ holds only on $\lambda$ and
$Y$ only on $\mu$, where $\lambda\ne\mu$.
Their composition notices the folds of
$\lambda$ and $\mu$ at $(\hub,v_\mu)$ and
$(v_\lambda,\hub)$, respectively.

Why can it not notice a third fold?
To notice $\lambda$ at $(\hub,v)$, $X$ must
hold on $\lambda$ and no other off-diagonal
color. Indeed, if $X[\nu]=1$ for some
$\nu\ne\lambda$ and $Y[\tau]=1$ for some
color $\tau$, the two-witness property
supplies a path
$\hub\xrightarrow{\nu}z\xrightarrow{\tau}v$
that survives the fold of $\lambda$.
Thus $z$ remains a witness for
$(X^{(\lambda)};Y^{(\lambda)})(\hub,v)$.
%
Likewise, noticing a distinct fold $\mu$
at $(u,\hub)$ requires $Y$ to hold only
on $\mu$. Hence noticing two distinct
folds forces $X$ and $Y$ to be supported
on $\{\lambda\}$ and $\{\mu\}$, respectively.
No third color occurs in either support,
and their supports are disjoint, ruling out
an additional fold noticed at $(\hub,\hub)$.
If $X$ and $Y$ share a color instead,
all possible noticed folds must coincide.
Thus $X;Y$ notices at most two folds.

\paragraph{Summary.}
Every one of the $2^k$ folds requires a noticing composition gate
(\S\ref{ov:types}), whereas each distinct composition gate notices at
most two (Lemma~\ref{lem:local}). Hence $2^k\le 2\gcost C$,
proving the lower bound of Theorem~\ref{thm:main} even with sharing.
% For these tests, the factor two is tight (Proposition~\ref{prop:promise}).
